% =========================================================================
% related_work.tex
% =========================================================================

\section{Related Work and Novelty Positioning}
\label{sec:related}

% -----------------------------------------------------------------
\subsection{KV Cache Quantization}

A substantial body of work addresses KV cache compression.
KIVI~\cite{liu2024kivi} applies asymmetric two-bit quantization with
per-channel tuning, achieving strong results on instruction-following
benchmarks.
KVQuant~\cite{hooper2024kvquant} targets 10M-token contexts via
sensitivity-aware per-channel quantization with nuq (non-uniform
quantization) codebooks.
AWQ~\cite{lin2024awq} focuses on activation-aware weight quantization
but its per-channel scaling insights extend to cache compression.
TurboQuant~\cite{zandieh2025turboquant} and
PolarQuant~\cite{han2025polarquant} advance this line by exploiting
rotation-induced distributional regularity for data-oblivious
compression.
All of the above validate on models with $\dhead \geq 128$.

% -----------------------------------------------------------------
\subsection{Community Findings on CLT Convergence at Small Head Dimensions}

Two community forks of TurboQuant have informally documented degraded
performance at $\dhead = 64$.
The AmesianX implementation~\cite{amesianx2025turboquant} for
\texttt{llama.cpp} includes an auto-fallback mechanism that reverts the
Key cache to \texttt{q8\_0} quantization when $\dhead = 64$, on the
grounds that CLT convergence is insufficient for the WHT rotation to
yield the expected Beta distribution.
The OnlyTerp fork~\cite{onlyterp2025turboquant} similarly documents
that the QJL residual correction introduces variance that outweighs
its bias-removal benefit at $\dhead \leq 64$.

These are non-peer-reviewed software observations, not systematic
empirical studies.
They do not propose a principled correction beyond falling back to
a different quantization format, and they do not test
cross-architecture combinations with Memory Caching.

% -----------------------------------------------------------------
\subsection{Novelty Claims}

I state precisely what is and is not novel in this work.

\paragraph{What is known.}
The observation that QJL residual correction degrades at small head
dimensions has been noted informally in community
implementations~\cite{amesianx2025turboquant, onlyterp2025turboquant}
and in discussions around the vLLM inference engine.
The general principle that CLT convergence weakens at lower
dimensionality is well established in probability theory.

\paragraph{Contributions.}

To my knowledge, no published work has proposed:
\begin{enumerate}[noitemsep,topsep=1pt,parsep=0pt,partopsep=0pt]
\item A systematic empirical characterization of TurboQuant,
      QJL, PolarQuant, and FibQuant failure modes specifically at
      $\dhead = 64$, with controlled comparisons across six
      quantization schemes.
\item An Adaptive QJL mechanism that auto-disables the residual
      correction based on head dimension, rather than falling back
      to a different quantization format entirely.
\item MC-TurboQuant: a combination of WHT-based vector quantization
      with GRM segmented caching from the Memory Caching framework.
\item An empirical demonstration that FibQuant-style angular codebooks
      collapse at $d = 64$, with quantitative comparison to scalar
      alternatives.
\end{enumerate}

All novelty claims are hedged: I state ``to my knowledge'' because
web searches, while thorough, cannot guarantee completeness.
If any of the above contributions have been independently proposed
elsewhere, I welcome corrections.
